\documentclass[11pt]{article}
\usepackage[margin=1in]{geometry}
\usepackage{times}
\usepackage{latexsym}
\usepackage{natbib}
\usepackage{hyperref}
\usepackage{url}
\usepackage{graphicx}
\usepackage{booktabs}
\usepackage{amsmath}
\usepackage{multirow}
\usepackage{xcolor}

% Highlight regressions in tables
\newcommand{\regress}[1]{\textcolor{red}{\textbf{#1}}}
\newcommand{\improve}[1]{\textcolor{black}{#1}}

\title{Valid JSON, Wrong Answer:\\Per-Role Regression Detection and\\Linguistic Presupposition Labeling for Structured Output}

\author{Breck Baldwin \\
  Independent Researcher \\
  New York, NY \\
  \texttt{breckbaldwin@gmail.com}}

\begin{document}
\maketitle

\begin{abstract}
Grammar-constrained decoding guarantees syntactically valid structured output from language models. Aggregate loss metrics show clear improvement after LoRA fine-tuning. But a considerable surprise hides in broad metric improvement---structured components don't necessarily share in the benefits of fine-tuning, and the result gets worse with scale.

We decompose fine-tuning loss by grammar role---structural tokens, schema keys, enumeration values, booleans, and free text---for Qwen~2.5 models fine-tuned to generate JSON at 0.5B, 7B, and 32B parameter scales. At 32B, we find that fine-tuning \emph{degrades} boolean prediction by 130\% on a flight booking schema while aggregate loss improves by 69\%. The model overfits to the majority boolean value in training data.

The regression emerges only at scale: at 0.5B, boolean performance improves 53\%; at 7B, it stalls; at 32B, it collapses. Lexical analysis reveals the root cause: 77\% of training examples force the model to predict a field value with no supporting evidence in the input. We propose \emph{presupposition labeling}---relabeling underspecified training targets as \texttt{ambiguous}---which completely eliminates the regression with no architectural changes.

We show that per-grammar-role decomposition is straightforward to compute post-hoc, reveals failure modes invisible to standard evaluation, and can motivate targeted mitigations. We release code, data, and the PyPI module \texttt{slotloss} to enable this analysis on any fine-tuned structured output model.
\end{abstract}

\section{Introduction}

Large language models are increasingly used to generate structured output: JSON for APIs, database queries, configuration files, clinical forms. Grammar-constrained decoding systems---Guidance \cite{guidance2024}, Outlines \cite{willard2023outlines}, SGLang \cite{zheng2024sglang}, XGrammar \cite{xgrammar2024}---can guarantee that every generated token is syntactically valid. This is a solved problem.

Fine-tuning with parameter-efficient methods such as LoRA \cite{hu2021lora} improves content quality, as measured by aggregate loss or accuracy metrics. The standard workflow is: fine-tune, observe aggregate loss decrease, confirm valid JSON via constrained decoding, deploy.

This workflow has a gap. Aggregate metrics mix \emph{trivially correct} decisions---structural tokens like \texttt{\{}, \texttt{\}}, \texttt{:}, \texttt{,} that have few valid options---with \emph{substantively uncertain} decisions like enum selection, boolean prediction, and free-text generation. A model can achieve large aggregate improvements by perfecting trivial structural predictions while degrading on the predictions that actually matter. Aggregate metrics will show improvement.

We demonstrate this empirically. Using the Schema-Guided Dialogue dataset \cite{rastogi2020sgd}, we fine-tune Qwen~2.5 Instruct models at three scales (0.5B, 7B, 32B) to generate JSON from dialogue context. We decompose per-token loss by grammar role and find:

\begin{itemize}
\item At 32B, fine-tuning degrades boolean prediction by 130\% on a flight booking schema (\texttt{refundable}: \texttt{True}/\texttt{False}), while aggregate loss improves by 69\%.
\item The regression is scale-dependent: at 0.5B, boolean performance improves 53\%; at 7B, it is marginal ($-$10\%); at 32B, it collapses.
\item Enumeration value prediction stalls at 32B ($-$8.6\%), approaching regression, while improving substantially at smaller scales.
\item A second schema (restaurant booking) shows no regression---the phenomenon is schema- and data-dependent, reinforcing the need for per-role analysis rather than aggregate evaluation.
\end{itemize}

Our decomposition method is post-hoc: parse the generated JSON, categorize each token by grammar role, compute per-role loss. It requires no changes to training and works with any model and schema. We release code and data.\footnote{Code and data: \texttt{https://github.com/breckbaldwin/strict-ft-eval}}

\section{Background}

\subsection{Grammar-Constrained Decoding}

Grammar-constrained decoding systems enforce structural validity during generation by masking tokens that would violate the target grammar. Systems including Guidance \cite{guidance2024}, Outlines \cite{willard2023outlines}, and SGLang \cite{zheng2024sglang} guarantee that output conforms to a JSON Schema, SQL grammar, or other formal specification.

These systems prevent \emph{syntactic} errors. They cannot prevent \emph{semantic} errors within the valid output space. A model constrained to valid JSON can produce \texttt{\{"refundable": "True"\}} when the correct answer is \texttt{\{"refundable": "False"\}}---both are valid JSON, both match the schema's constraint, but one is wrong. 

\subsection{LoRA Fine-Tuning for Structured Output}

Low-Rank Adaptation \cite{hu2021lora} is a standard method for fine-tuning large language models on domain-specific tasks. For structured output, standard practice is: apply LoRA to attention projections, train with cross-entropy loss on target sequences, and evaluate with aggregate metrics (mean loss, exact match, ROUGE).

The implicit assumption is that if aggregate loss improves, all aspects of output quality improve. We show this assumption is violated.

\subsection{JSON as Structured Output}

A JSON document against a fixed schema contains tokens playing distinct grammar roles:

\begin{itemize}
\item \textbf{Structural}: \texttt{\{}, \texttt{\}}, \texttt{[}, \texttt{]}, \texttt{:}, \texttt{,}---near-deterministic, few valid options
\item \textbf{Quote}: string delimiters---deterministic
\item \textbf{Key}: object key characters---constrained to schema-defined keys
\item \textbf{Enum value}: categorical values from a fixed set
\item \textbf{Boolean}: \texttt{True} or \texttt{False} strings
\item \textbf{Free text}: unconstrained string content (names, addresses)
\item \textbf{Whitespace}: formatting---trivial
\end{itemize}

These roles differ enormously in difficulty. Structural tokens and whitespace are near-deterministic and contribute many tokens to the sequence. Free text and boolean/enum values carry genuine uncertainty but contribute fewer tokens. Aggregate loss can be dominated by the high-token-count trivial roles or not. Our observations are very schema/data dependent.

\section{Method: Per-Grammar-Role Decomposition}

\subsection{Grammar Role Assignment}

Given a JSON string and its schema, we assign each \emph{character} a grammar role by walking the string with a recursive-descent parser that tracks:

\begin{enumerate}
\item Whether the current position is inside a key or value
\item For values: whether the enclosing key maps to an enum, boolean, or free-text field in the schema
\item Structural characters, quotes, and whitespace are assigned directly
\end{enumerate}

\subsection{Mapping to Tokens}

We tokenize the JSON with the model's standard BPE tokenizer. For each token, we determine the character span it covers and assign the grammar role of the token's first character. Tokens spanning multiple grammar roles (e.g., the BPE token \texttt{'":"'} spanning quote + colon + quote) are assigned to the first role.

This is a known limitation of post-hoc analysis: cross-terminal tokens prevent exact per-role attribution. We note that such tokens are a small fraction of the total and that the dominant roles (KEY, ENUM, FREE\_TEXT) rarely share tokens with other roles. 

Note that \textbf{Boolean} for this experiment had a 1-1 mapping from BPE tokenization to JSON terminal across all three Qwen models so the result is not an artifact of surface tokenization differences. 

\subsection{Per-Role Loss Computation}

For each example, we compute teacher-forced per-token cross-entropy loss on the target JSON (prompt tokens are masked). We aggregate loss by grammar role:

\begin{equation}
\mathcal{L}_r = \frac{1}{|T_r|} \sum_{t \in T_r} \ell(t)
\end{equation}

where $T_r$ is the set of tokens assigned to grammar role $r$ and $\ell(t)$ is the per-token cross-entropy loss.

\subsection{Applicability}

This method works with any fine-tuned model, any JSON schema, and any tokenizer. It requires no modifications to training. The only inputs are: the model, a prompt, the target JSON, and the schema definition. The analysis adds negligible overhead---it is a single forward pass with token-level loss recording.

\section{Experimental Setup}

\subsection{Data}

We use the Schema-Guided Dialogue (SGD) dataset \cite{rastogi2020sgd}, extracting two schemas:

\textbf{Restaurants\_1} (11 keys): 3 enum fields (cuisine, price\_range, party\_size), 2 boolean fields (serves\_alcohol, has\_live\_music), 6 free-text fields. An intuitive domain with moderate schema complexity.

\textbf{Flights\_1} (16 keys): 4 enum fields (passengers, seating\_class, number\_stops, airlines), 1 boolean field (refundable), 11 free-text fields. A complex domain with a high-cardinality enum (8 airlines).

For each schema, we extract 250 training and 50 test examples. Each example pairs a dialogue context with a service result JSON object.

\subsection{Models and Training}

We use Qwen~2.5 Instruct at three scales: 0.5B, 7B, and 32B parameters. Fine-tuning uses standard LoRA via PEFT \cite{peft2023}: rank 16, alpha 32, targeting \texttt{q\_proj} and \texttt{v\_proj}, trained for 10 epochs with learning rate $10^{-4}$. No custom adapter architectures or training modifications.

Grammar-constrained decoding uses llguidance \cite{guidance2024} at inference time. All models produce 100\% valid JSON on held-out data.

\section{Results}

\subsection{Aggregate Metrics Tell a Positive Story}

Table~\ref{tab:aggregate} shows aggregate loss improvement across all conditions. Fine-tuning produces substantial improvements at every scale: 59--81\% reduction in mean loss. All fine-tuned models produce valid JSON. Standard evaluation says: fine-tuning works.

\begin{table}[t]
\centering
\small
\begin{tabular}{llccc}
\toprule
\textbf{Schema} & \textbf{Scale} & \textbf{Base} & \textbf{FT} & \textbf{Change} \\
\midrule
\multirow{3}{*}{Rest.} & 0.5B & 0.970 & 0.394 & $-$59\% \\
                        & 7B   & 1.090 & 0.359 & $-$67\% \\
                        & 32B  & 0.856 & 0.370 & $-$57\% \\
\midrule
\multirow{3}{*}{Flights} & 0.5B & 0.733 & 0.141 & $-$81\% \\
                          & 7B   & 0.645 & 0.161 & $-$75\% \\
                          & 32B  & 0.554 & 0.174 & $-$69\% \\
\bottomrule
\end{tabular}
\caption{Aggregate mean loss (all grammar roles). Fine-tuning improves all conditions. Standard evaluation sees no problems.}
\label{tab:aggregate}
\end{table}

\subsection{Per-Grammar-Role Decomposition Reveals a Hidden Regression}

Table~\ref{tab:flights_decomp} decomposes loss by grammar role for the Flights schema. At 32B, boolean prediction loss \emph{increases} from 0.457 to 1.050---a 130\% degradation---while aggregate loss decreases 69\%.

\begin{table}[t]
\centering
\small
\begin{tabular}{lcccccc}
\toprule
 & \multicolumn{2}{c}{\textbf{0.5B}} & \multicolumn{2}{c}{\textbf{7B}} & \multicolumn{2}{c}{\textbf{32B}} \\
\cmidrule(lr){2-3} \cmidrule(lr){4-5} \cmidrule(lr){6-7}
\textbf{Role} & Base & FT & Base & FT & Base & FT \\
\midrule
STRUCT  & 4.84 & 0.00 & 6.08 & 0.00 & 5.33 & 0.00 \\
QUOTE   & 0.13 & 0.00 & 0.16 & 0.00 & 0.10 & 0.00 \\
KEY     & 1.17 & 0.00 & 0.77 & 0.00 & 0.47 & 0.00 \\
ENUM    & 0.48 & 0.26 & 0.55 & 0.34 & 0.33 & 0.30 \\
\textbf{BOOL} & \textbf{0.80} & \textbf{0.37} & \textbf{0.43} & \textbf{0.39} & \textbf{0.46} & \regress{1.05} \\
FREE    & 1.31 & 0.52 & 1.26 & 0.59 & 1.33 & 0.63 \\
WS      & 0.04 & 0.00 & 0.04 & 0.00 & 0.02 & 0.00 \\
\midrule
TOTAL   & 0.73 & 0.14 & 0.64 & 0.16 & 0.55 & 0.17 \\
\bottomrule
\end{tabular}
\caption{Flights: per-grammar-role loss at three scales. At 32B, boolean loss \regress{increases 130\%} while aggregate loss decreases 69\%. The regression is invisible without decomposition.}
\label{tab:flights_decomp}
\end{table}

The Flights schema has one boolean field: \texttt{refundable} (\texttt{True}/\texttt{False}). At 32B, the pretrained model predicts this field reasonably well (0.457 loss). Fine-tuning disrupts this: the model most likely overfits to the majority boolean value in the training data, producing worse predictions than the untrained baseline. 

\subsection{The Regression Is Scale-Dependent}

Figure~\ref{fig:scaling} shows boolean prediction change across model scales on the Flights schema:

\begin{table}[h]
\centering
\small
\begin{tabular}{lccc}
\toprule
\textbf{Scale} & \textbf{Base} & \textbf{FT} & \textbf{Change} \\
\midrule
0.5B & 0.796 & 0.371 & $-$53\% \\
7B   & 0.432 & 0.389 & $-$10\% \\
32B  & 0.457 & \regress{1.050} & \regress{+130\%} \\
\bottomrule
\end{tabular}
\caption{Boolean (refundable) loss on Flights across scales. At 0.5B, fine-tuning helps substantially. At 7B, it barely helps. At 32B, it causes a severe regression. Larger models have more to lose.}
\label{fig:scaling}
\end{table}

At 0.5B, the base model is poor at boolean prediction (0.796 loss) and fine-tuning helps substantially ($-$53\%). At 7B, the base model is better (0.432) and fine-tuning barely helps ($-$10\%). At 32B, fine-tuning \emph{destroys} the base model's boolean competency.

In more detail, training data contains 76\% False, 24\% True for the refundable field. The fine-tuned 32B model's boolean loss (1.05) is consistent with assigning approximately 2\% probability to True — far more skewed than the training distribution itself. Fine-tuning amplified the majority-class bias rather than learning the conditional distribution.

The pattern is intuitive: larger pretrained models have stronger existing competencies. Fine-tuning on a small dataset (250 examples) can override these competencies with memorized patterns. The better the base model already is, the more fine-tuning has to disrupt.

\subsection{Enumeration Values Are Approaching Regression}

A second warning sign appears in the ENUM\_VALUE role on Flights:

\begin{table}[h]
\centering
\small
\begin{tabular}{lccc}
\toprule
\textbf{Scale} & \textbf{Base} & \textbf{FT} & \textbf{Change} \\
\midrule
0.5B & 0.478 & 0.259 & $-$46\% \\
7B   & 0.554 & 0.337 & $-$39\% \\
32B  & 0.331 & 0.303 & $-$9\% \\
\bottomrule
\end{tabular}
\caption{Enum value loss on Flights. Improvement diminishes with scale, approaching regression at 32B.}
\label{tab:enum_trend}
\end{table}

The improvement from fine-tuning decreases monotonically with scale: $-$46\% at 0.5B, $-$39\% at 7B, $-$9\% at 32B. The Flights schema has high-cardinality enums (8 airlines, 4 seating classes). At 32B, the pretrained model's enum predictions are already reasonable (0.331 loss), and fine-tuning adds negligible value. With more training epochs or slightly different data, this could tip into regression.

\subsection{Not All Schemas Show Regression}

Table~\ref{tab:restaurants_decomp} shows the Restaurants schema, where \emph{no role regresses} at any scale.

\begin{table}[t]
\centering
\small
\begin{tabular}{lcccccc}
\toprule
 & \multicolumn{2}{c}{\textbf{0.5B}} & \multicolumn{2}{c}{\textbf{7B}} & \multicolumn{2}{c}{\textbf{32B}} \\
\cmidrule(lr){2-3} \cmidrule(lr){4-5} \cmidrule(lr){6-7}
\textbf{Role} & Base & FT & Base & FT & Base & FT \\
\midrule
STRUCT  & 4.94 & 0.00 & 6.25 & 0.00 & 5.52 & 0.00 \\
QUOTE   & 0.22 & 0.04 & 0.33 & 0.04 & 0.17 & 0.02 \\
KEY     & 1.21 & 0.00 & 1.32 & 0.00 & 0.67 & 0.00 \\
ENUM    & 1.73 & 0.36 & 1.70 & 0.56 & 1.59 & 0.51 \\
BOOL    & 0.81 & 0.36 & 1.24 & 0.63 & 0.74 & 0.50 \\
FREE    & 1.76 & 1.21 & 1.86 & 1.05 & 1.71 & 1.11 \\
WS      & 0.02 & 0.00 & 0.08 & 0.00 & 0.03 & 0.00 \\
\midrule
TOTAL   & 0.97 & 0.39 & 1.09 & 0.36 & 0.86 & 0.37 \\
\bottomrule
\end{tabular}
\caption{Restaurants: per-grammar-role loss at three scales. No role regresses. The boolean regression is schema- and data-dependent.}
\label{tab:restaurants_decomp}
\end{table}

Restaurants has two boolean fields (serves\_alcohol, has\_live\_music) compared to Flights' one (refundable). The presence or absence of regression depends on the schema, the training data distribution, and the base model's prior competency on each field. This is precisely why per-grammar-role decomposition is necessary: you cannot predict \emph{which} role will regress without measuring it.

\subsection{Cross-Domain Validation: Contract Clause Extraction (CUAD)}\label{sec:cuad}

To test whether the per-role phenomena generalise beyond task-oriented dialogue, we apply the decomposition method to the Contract Understanding Atticus Dataset \cite{hendrycks2021cuad} -- a corpus of 510 commercial contracts with human-annotated clause-level labels. We construct an 11-field JSON schema: eight boolean clause-presence indicators spanning 4\% to 78\% Yes-rate, and three normalised enum fields (\texttt{governing\_law}, 12 values; \texttt{renewal\_term}, 6; \texttt{expiration\_type}, 3). The enum fields include a \texttt{not\_specified} value -- the natural analog of the \texttt{ambiguous} label introduced for presupposition labelling (Section~\ref{sec:pcl}). The task is document-level: given a contract excerpt, output the JSON. Train/test split is 144/49, stratified by total True-count on booleans; input shape is considerably longer ($\sim$2k tokens) than the dialogue schemas ($\sim$0.5k tokens).

\textbf{Per-field baseline accuracy at 0.5B ranges from 2\% to 94\%.} Aggregate accuracy on the 11-field schema would average these to an uninformative middle value; the aggregate metric obscures nearly every field-level story. This reproduces the paper's central thesis on a document-level input, demonstrating that the pattern is not an artefact of dialogue state tracking. Crucially, aggregate accuracy at 7B (65.1\%) and 32B (66.4\%) is nearly identical -- a developer using aggregate metrics would conclude 7B is sufficient. Per-field analysis contradicts this conclusion.

\textbf{Boolean fields partition into two regimes by the sign of the pretrained prior.} For fields where the base model's default (False) matches the majority gold class -- e.g., \texttt{has\_most\_favored\_nation} with 4\% True gold -- baseline accuracy is 94\% across all scales and per-field diagnostics confirm strong calibration. For fields where the base model's default contradicts the majority gold class -- e.g., \texttt{has\_anti\_assignment} with 78\% True gold -- baseline accuracy is 31\% at 0.5B and \emph{further regresses to 24\% at 7B} before partially recovering to 35\% at 32B. The 7B dip is accompanied by a mean margin of 0.99: the model becomes \emph{confidently more wrong} at intermediate scale before partially recovering. This non-monotonic failure is invisible to aggregate accuracy, which rises monotonically with scale on this schema.

\begin{table}[h]
\centering
\small
\begin{tabular}{lccc}
\toprule
\textbf{Field} & \textbf{0.5B} & \textbf{7B} & \textbf{32B} \\
\midrule
\texttt{has\_most\_favored\_nation}  & 94\% (0.38) & 96\% (0.90) & 94\% (0.94) \\
\texttt{has\_liquidated\_damages}    & 86\% (0.51) & 90\% (0.75) & 90\% (0.83) \\
\texttt{has\_exclusivity}            & 65\% (0.23) & 82\% (0.76) & 92\% (0.78) \\
\texttt{has\_anti\_assignment}       & 31\% (0.20) & \regress{24\% (0.99)} & 35\% (0.96) \\
\texttt{governing\_law} (12-way)     & 14\% (0.21) & 20\% (0.57) & 24\% (0.74) \\
\texttt{renewal\_term} (6-way)       & \regress{2\% (0.45)} & 51\% (0.65) & 55\% (0.72) \\
\texttt{expiration\_type} (3-way)    & 49\% (0.26) & 67\% (0.67) & 57\% (0.80) \\
\bottomrule
\end{tabular}
\caption{CUAD baseline accuracy (mean margin) across scales for selected fields. Regressions in red. \texttt{has\_anti\_assignment} shows per-field regression from 0.5B to 7B invisible to the $\sim$65\% aggregate at both scales; \texttt{renewal\_term} shows a scale-dependent acquisition of the \texttt{not\_specified} sentinel between 0.5B and 7B.}
\label{tab:cuad_per_field_scale}
\end{table}

\subsection{Why Aggregate Metrics Hide the Regression}

The mechanism is simple. In the Flights 32B case:
\begin{itemize}
\item \textbf{STRUCTURAL} tokens: 100 tokens, loss drops from 5.33 to 0.00 ($-$100\%)
\item \textbf{WHITESPACE}: 1,950 tokens, loss drops from 0.02 to 0.00 ($-$100\%)
\item \textbf{BOOLEAN}: 50 tokens, loss increases from 0.46 to 1.05 (+130\%)
\end{itemize}

The massive improvements on high-token-count trivial roles (structural, whitespace) swamp the regression on the low-token-count substantive role (boolean). The aggregate is a weighted average where the weights are token counts, and trivial roles contribute far more tokens than substantive ones.

\section{Discussion}

\subsection{Why the Boolean Regresses: Lexical Analysis}

To understand the regression mechanism, we examined the training data for lexical support of the \texttt{refundable} field. The training set contains 76\% False and 24\% True values. We categorized each example by whether the conversation contains any discussion of refundability:

\begin{table}[h]
\centering
\small
\begin{tabular}{lcc}
\toprule
\textbf{Context} & \textbf{True} & \textbf{False} \\
\midrule
Refundability discussed (57 examples) & 47\% & 53\% \\
Not discussed (193 examples) & 17\% & 83\% \\
\bottomrule
\end{tabular}
\caption{Boolean value distribution conditioned on lexical support. When refundability is discussed, the distribution is roughly balanced. When not discussed, the field is a property of the API result with no conversational cue.}
\label{tab:lexical}
\end{table}

When users explicitly discuss refundability (``I need refundable tickets''), the True/False distribution is roughly 50/50 and the model has lexical cues to predict from. But in 77\% of examples (193/250), refundability is never mentioned --- the \texttt{refundable} field is simply a property of the flight returned by the API, with no signal in the conversation to predict it from.

This means 54\% of all True cases (32/59) have \emph{no lexical support} in the input. The ``correct'' prediction for these cases is genuine uncertainty --- approximately $-\log(0.5) = 0.69$ loss if the model honestly represents its ignorance. Instead, the fine-tuned 32B model assigns approximately 2\% probability to True (consistent with observed loss of 1.05), far more extreme than either the training distribution (24\% True) or the conversational evidence.

Fine-tuning does not merely memorize the training distribution --- it \emph{amplifies} the majority-class bias. The gradient signal from 193 undiscussable examples overwhelms the 57 examples with genuine lexical cues, pushing the model toward false confidence on an inherently uncertain field.

\subsection{Presupposition Labeling: A Mitigation}\label{sec:pcl}

The lexical analysis suggests a mitigation: if the input does not provide evidence for a field's value, the training label should reflect this rather than forcing a commitment. We draw on the \emph{uniqueness presupposition} from pronoun resolution \cite{baldwin1997cogniac}: a singular pronoun presupposes exactly one valid referent; when multiple referents are equally valid, the system should abstain rather than guess.

Applied to structured output: for each constrained field, we ask whether the input provides sufficient evidence to determine the value. If not, we relabel the training target as \texttt{ambiguous} --- a third option alongside the original values. For the \texttt{refundable} field, the 57 examples where refundability is discussed retain their True/False labels; the 193 examples where it is not discussed are relabeled to \texttt{ambiguous}. The schema is updated accordingly: \texttt{refundable: ["True", "False", "ambiguous"]}.

We call this \textbf{presupposition labeling}: expanding the label set based on whether the input satisfies the evidential presupposition for a field's value.

\begin{table}[h]
\centering
\small
\begin{tabular}{lccc}
\toprule
\textbf{Scale} & \textbf{2-way Std LoRA} & \textbf{3-way Std LoRA} & \textbf{Change} \\
\midrule
0.5B & 0.26 & 0.21 & $-$19\% \\
7B   & 0.34 & 0.27 & $-$21\% \\
32B  & \regress{1.05 (+130\%)} & \improve{0.19 (no regression)} & \textbf{eliminated} \\
\bottomrule
\end{tabular}
\caption{ENUM/BOOLEAN loss with and without presupposition labeling. At 32B, the +130\% regression is completely eliminated. No architectural changes --- only the training labels are modified.}
\label{tab:pcl}
\end{table}

Table~\ref{tab:pcl} shows the result: presupposition labeling \textbf{completely eliminates the boolean regression} at 32B. Standard LoRA with 3-way labels achieves 0.19 ENUM loss --- better than any 2-way condition --- with no regression at any scale. The model learns to output \texttt{ambiguous} for undiscussed cases instead of memorizing the majority class.

Critically, this requires \emph{no changes to the model architecture or training procedure}. Only the training labels are modified. The regression was not caused by the training method --- it was caused by training labels that forced commitment on underspecified inputs. Presupposition labeling is a standalone data preparation technique compatible with any fine-tuning method.

This does not diminish the value of per-grammar-role decomposition: without the decomposition, the regression would not have been detected, and without the lexical analysis it motivated, the labeling problem would not have been identified. The diagnostic method remains essential even when the treatment is simple.

\textbf{Extending to CUAD reveals field-specific variance.} Applying PCL fine-tuning to the 11-field CUAD schema (Section~\ref{sec:cuad}) across the same three scales produces a mixed picture. Aggregate accuracy rises modestly (+12/+3/+5 pp at 0.5B/7B/32B, vs. +30/+19/+20 on Flights), and the per-field breakdown explains why: PCL-FT's benefit is concentrated on fields whose pretrained prior conflicts with the task distribution (\texttt{has\_anti\_assignment} 35$\rightarrow$61\% at 32B; \texttt{governing\_law} 25$\rightarrow$39\% at 32B). On fields whose pretrained behaviour was already correct (\texttt{has\_most\_favored\_nation}, \texttt{has\_liquidated\_damages}), PCL-FT is neutral. On some fields PCL-FT \emph{regresses} (\texttt{expiration\_type} falls from 49/67/57\% baseline to 43/55/39\% after PCL-FT at 0.5B/7B/32B; \texttt{has\_non\_compete} 80/88/82\% baseline drops to 65/76/65\% PCL-FT). The regressions concentrate on fields whose training signal is weakened by truncation: contracts are truncated to the first 8000 characters for training efficiency, and the diagnostic clauses for these fields often sit deeper in the document. We confirm this hypothesis directly: re-evaluating the same 32B PCL-FT model on full-context CUAD test contracts lifts aggregate accuracy from 71\% to 85\% and substantially recovers both regressions (\texttt{expiration\_type} 39$\rightarrow$53\%, \texttt{has\_non\_compete} 65$\rightarrow$79\% at 32B). The training-time truncation is the bottleneck, not PCL-FT itself. The finding is important for the deployment posture: PCL-FT is not a uniform win, per-field evaluation remains necessary even after intervention, and training-time truncation choices affect which fields PCL-FT can rescue.

\subsection{Presupposition Labeling as a Third Category}

Presupposition labeling is neither supervised labeling nor an autoregressive technique. It occupies a distinct category:

\begin{itemize}
\item \textbf{Autoregressive training}: the signal comes from predicting the next token given preceding context.
\item \textbf{Supervised labeling}: the signal comes from human annotator judgment on individual examples.
\item \textbf{Presupposition labeling}: the signal comes from the structural relationship between the schema and the input. The schema defines what values are valid; the presupposition check determines whether the input \emph{licenses} commitment to any specific value. No human judges individual examples --- a rule derived from the schema-input pair is applied uniformly.
\end{itemize}

The JSON schema already operates outside the autoregressive framework: grammar-constrained decoding restricts what tokens the model may produce. Presupposition labeling extends this principle to training: the schema restricts not only what values are \emph{structurally valid} but which training examples provide \emph{evidential license} to commit to a specific value. This is a schema-mediated evidential constraint on the training signal.

In our experiment, the presupposition check is implemented as a lexical cue (does the conversation mention ``refund''?). This is a crude proxy, but the principle is general: for any constrained field, one can ask whether the input contains evidence distinguishing among the field's allowed values. More sophisticated implementations could derive presupposition checks automatically from the schema --- for example, by measuring mutual information between input features and field values, and relabeling examples where the mutual information is below a threshold.

The connection to linguistic theory is direct. CogNIAC's uniqueness presupposition \cite{baldwin1997cogniac} treats ambiguity as a semantic violation, not uncertainty to be resolved: if multiple referents are equally valid for a singular pronoun, the system halts. Presupposition labeling applies the same principle to structured output: if the input does not determine a field's value, the training signal should encode abstention (\texttt{ambiguous}), not force a commitment that the model can only satisfy by memorizing the majority class.

\subsection{Margin-Based Evidential Gating at Inference Time}\label{sec:gating}

Presupposition labeling eliminates the regression at training time. A complementary question is whether the evidential signal is already present in the pretrained output distribution -- and can be exploited at inference without fine-tuning.

At each constrained-field position, we measure the margin $m = p_1 - p_2$, where $p_1$ and $p_2$ are the top and second-highest probabilities over the field's allowed values. Low margin indicates the model is indifferent between candidates; high margin indicates a confident commitment. We commit to the top value when $m \ge \theta$ and abstain otherwise.

\begin{table}[h]
\centering
\small
\begin{tabular}{lccc}
\toprule
 & \textbf{Gold = True/False} & \textbf{Gold = ambiguous} & \\
\textbf{Group} & mean margin & mean margin & baseline acc \\
\midrule
Discussed (Flights, $n=14$) & 0.45 & --- & 60\% \\
Undiscussed (Flights, $n=36$) & --- & 0.20 & 0\% (forced commit) \\
\bottomrule
\end{tabular}
\caption{Base Qwen 2.5 0.5B, Flights\_1 \texttt{refundable}. The margin doubles when the input contains lexical evidence (discussed); it stays flat when the field is underspecified. The signal for evidential gating is present in the base model without any fine-tuning.}
\label{tab:gating_flights}
\end{table}

On Flights\_1, the margin discriminates discussed from undiscussed examples cleanly (Table~\ref{tab:gating_flights}). At $\theta = 0.35$, the base model correctly abstains on 100\% of undiscussed gold examples while retaining 71\% committed accuracy on the discussed subset (up from 43\% baseline under forced commitment).

On CUAD's \texttt{expiration\_type} field the same operating point achieves 76\% abstention on \texttt{not\_specified} golds and 82\% committed accuracy on the specified subset -- a replication of the Flights\_1 result on a document-level task.

Not every field supports this strategy, however. On CUAD's \texttt{renewal\_term} field, where 66\% of gold labels are \texttt{not\_specified}, the base model's abstention rate at $\theta = 0.30$ varies sharply with scale: at 0.5B the pretrained distribution places near-zero probability on the \texttt{not\_specified} token at the schema-value position, yielding 0\% abstention on the 34 undiscussed examples and 2\% overall accuracy. By 7B, the same field reaches 53\% baseline accuracy, and by 32B, 55\%. The model learns to emit the sentinel value somewhere between 0.5B and 7B; further scale adds little. Adding a sentinel value to the schema is thus a scale-sensitive intervention: at small model sizes it is insufficient without fine-tuning; at 7B or above it becomes usable but does not fully close the gap (55\% vs. the 66\% rate of \texttt{not\_specified} in gold). When label-only approaches are insufficient, presupposition labeling at training time is the remaining lever.

\textbf{Margin calibration recovers at 32B.} A further observation from scaling margin gating across model sizes: the 7B scale produces the \emph{highest} mean margins on most constrained fields (often $> 0.95$), while 32B margins are systematically lower (typically $0.7$--$0.9$). Table~\ref{tab:margin_scaling} shows the pattern. Higher 7B margins reflect over-commitment: the model is confidently wrong on undiscussed cases (e.g., \texttt{has\_anti\_assignment} mean margin 0.99 at 22\% accuracy). At 32B the distribution over allowed values is less saturated, preserving evidential calibration that margin gating can read. This directly supports the proposition that 32B is more useful for inference-time gating than 7B despite their near-identical aggregate accuracies.

\begin{table}[h]
\centering
\small
\begin{tabular}{lccc}
\toprule
\textbf{Field} & \textbf{0.5B} & \textbf{7B} & \textbf{32B} \\
\midrule
\texttt{refundable} (Flights)                    & 0.25 & \regress{0.80} & 0.45 \\
\texttt{has\_anti\_assignment} (CUAD)            & 0.19 & \regress{0.99} & 0.93 \\
\texttt{has\_cap\_on\_liability} (CUAD)          & 0.24 & \regress{0.97} & 0.92 \\
\texttt{has\_most\_favored\_nation} (CUAD)       & 0.36 & 0.90 & 0.93 \\
\texttt{expiration\_type} (CUAD)                 & 0.25 & 0.68 & 0.76 \\
\bottomrule
\end{tabular}
\caption{Mean margin ($p_1 - p_2$) across scales. Peaks at 7B for most fields (red), indicating over-commitment. 32B margins are systematically lower and track evidence more faithfully, making margin-based abstention more informative at scale.}
\label{tab:margin_scaling}
\end{table}

Margin gating and presupposition labeling address complementary halves of the same underlying problem. The former reads the model's existing evidential calibration; the latter aligns the training signal with that calibration. Neither requires architectural changes. The two combine: margin gating works best at scales where calibration is preserved (32B and above), while presupposition labeling rescues smaller models and sharpens larger ones.

\subsection{Head-to-Head: Training-Time and Inference-Time Interventions}\label{sec:head_to_head}

We evaluated all four combinations of \{baseline, PCL-FT\} $\times$ \{argmax, margin gating\} on Flights, Restaurants, and CUAD at 0.5B / 7B / 32B. Table~\ref{tab:head_to_head} reports aggregate accuracy across every constrained field per dataset; margin gating uses $\theta = 0.30$ with coverage (fraction of fields committed on) reported alongside committed accuracy.

\begin{table}[t]
\centering
\small
\begin{tabular}{llcccc}
\toprule
 &  & \multicolumn{2}{c}{\textbf{Baseline}} & \multicolumn{2}{c}{\textbf{PCL-FT}} \\
\cmidrule(lr){3-4} \cmidrule(lr){5-6}
\textbf{Dataset} & \textbf{Scale} & argmax & + IT-1 ($\theta{=}0.30$) & argmax & + IT-1 ($\theta{=}0.30$) \\
\midrule
Flights     & 0.5B & 58\% & 74\% @68\%     & 88\% & \textbf{89\%} @97\% \\
Flights     & 7B   & 71\% & 72\% @96\%     & 90\% & \textbf{93\%} @94\% \\
Flights     & 32B  & 72\% & 84\% @82\%     & 92\% & \textbf{95\%} @95\% \\
\midrule
Restaurants & 0.5B & 31\% & 39\% @58\%     & 81\% & \textbf{87\%} @92\% \\
Restaurants & 7B   & 45\% & 48\% @93\%     & 83\% & \textbf{85\%} @96\% \\
Restaurants & 32B  & 52\% & 52\% @92\%     & 82\% & \textbf{83\%} @97\% \\
\midrule
CUAD        & 0.5B & 56\% & 64\% @58\%     & 68\% & \textbf{71\%} @77\% \\
CUAD        & 7B   & 65\% & 67\% @89\%     & 68\% & \textbf{70\%} @91\% \\
CUAD        & 32B  & 66\% & 67\% @92\%     & 71\% & \textbf{73\%} @93\% \\
\bottomrule
\end{tabular}
\caption{Aggregate accuracy across all constrained fields, per dataset $\times$ scale. Cells are argmax accuracy for argmax columns and ``committed accuracy @ coverage'' for the gated columns. PCL-FT + inference-time gating is the best cell in every row. CUAD here uses the 8000-character truncated prompt; full-context CUAD reaches 85\%/87\%@96\% at 32B (Section~\ref{sec:cuad}).}
\label{tab:head_to_head}
\end{table}

Three patterns are visible:

\begin{itemize}
\item \textbf{PCL-FT always helps.} Every row shows PCL-FT argmax $\ge$ baseline argmax. On Flights the gap is +19/+20 pp at 7B/32B, driven by \texttt{refundable} recovering from the 32B regression. Restaurants shows the largest gap (+38/+30 pp), driven by both boolean fields (\texttt{serves\_alcohol}, \texttt{has\_live\_music}) where ``alcohol'' and ``music'' carry unique evidential signal. CUAD's gap is the smallest (+3/+5 pp), partly because the truncated prompt cuts off cue evidence for several fields and partly because CUAD's existing \texttt{not\_specified} sentinel already absorbs much of the under-specification signal that PCL adds for the other two datasets.
\item \textbf{Margin gating always adds a few percentage points on top of PCL-FT.} Across all rows where both columns exist, gating adds 1--6 pp of committed accuracy for 3--8 pp of dropped coverage. The two interventions are not redundant: PCL-FT shifts the model to use the \texttt{ambiguous} sentinel during commitment; gating reads margin distance to detect the residual underconfident commitments that the model still emits.
\item \textbf{Gating helps most where commitment is least informative.} The two largest gating-only lifts are 0.5B Flights (58$\rightarrow$74\% at 68\% coverage) and 32B Flights (72$\rightarrow$84\% at 82\% coverage), both without any fine-tuning. The first is the small-model regime where fine-tuning is impractical; the second is the regime where evidential calibration is best preserved (Section~\ref{sec:gating}). In both, gating reads the model's existing uncertainty directly.
\end{itemize}

The combined intervention (PCL-FT + margin gating) reaches 70--95\% committed accuracy across our scales and datasets, dominating every other configuration. The $\theta = 0.30$ operating point is conservative; tightening it further trades coverage for precision.

\subsection{Remaining Open Directions}

Presupposition labeling addresses regressions caused by underspecified inputs. Other sources of per-role regression --- such as gradient domination by trivial structural tokens, or overfitting on small enum sets --- may require different interventions. Per-role learning rate schedules, per-role early stopping, or training-time structural awareness are potential directions. Addressing per-grammar-role regressions during the training process itself, rather than through data relabeling alone, remains an open problem.

\subsection{Implications for Deployment}

Any structured output task where aggregate metrics mix trivially-correct decisions with substantively-uncertain decisions is vulnerable to this failure mode. The harder decisions can regress while aggregate metrics improve.

We suggest per-grammar-role thresholds as deployment criteria:
\begin{itemize}
\item Deploy when STRUCTURAL $< 0.01$ and KEY $< 0.1$
\item Flag for review when ENUM\_VALUE $> 0.3$
\item Require human verification when BOOLEAN $> 0.5$
\end{itemize}

Such criteria are more informative than a single aggregate threshold and can be computed with one additional forward pass per evaluation example.

\subsection{The Medical Analogy}

Aggregate loss is a clean bill of health. Per-grammar-role decomposition is the blood panel. A patient---or a model---can appear healthy on aggregate while a specific system is failing. The decomposition costs little and can prevent deploying a model with a critical blind spot.

\section{Related Work}

\textbf{Grammar-constrained decoding.} Guidance \cite{guidance2024}, Outlines \cite{willard2023outlines}, SGLang \cite{zheng2024sglang}, and XGrammar \cite{xgrammar2024} guarantee syntactic validity at decode time. Our work is complementary: we evaluate \emph{content quality} decomposed by grammar role, which constrained decoding cannot assess.

\textbf{LoRA and parameter-efficient fine-tuning.} LoRA \cite{hu2021lora} and its variants (QLoRA \cite{dettmers2023qlora}, DoRA \cite{liu2024dora}) are standard for fine-tuning large models. We use unmodified LoRA and show that per-role regressions arise from the standard training procedure, not from any specific adapter variant.

\textbf{Evaluation methodology for LLMs.} Growing work addresses pitfalls in LLM evaluation: benchmark contamination \cite{jacovi2023contamination}, metric gaming, and the gap between aggregate scores and actual capability. Our contribution is specific to structured output: aggregate loss hides per-grammar-role regressions.

\textbf{Per-class analysis in classification.} Decomposing performance by class is standard in classification (confusion matrices, per-class F1). Our work applies this principle to structured generation, where ``classes'' are grammar roles rather than output labels.

\textbf{Calibration and confidence estimation.} Work on LLM calibration \cite{kadavath2022calibration} examines whether model confidence matches accuracy. Our per-grammar-role decomposition is related: it reveals where the model's predictions are reliable (structural tokens) versus unreliable (content tokens), at a finer granularity than aggregate calibration.

\section{Conclusion}

Standard evaluation of fine-tuned structured output models hides per-grammar-role regressions behind aggregate metrics. We demonstrate that a 32B model fine-tuned for JSON extraction degrades boolean prediction by 130\% while appearing to improve by 69\% overall. The regression emerges at scale, is schema-dependent, and is completely invisible to aggregate evaluation.

A simple post-hoc decomposition---parsing output by grammar role and computing per-role loss---reveals the regression and costs one forward pass. Lexical analysis of the training data reveals the root cause: 77\% of training examples force the model to predict a field value with no supporting evidence in the input.

We show that \emph{presupposition labeling}---relabeling underspecified training targets as \texttt{ambiguous}---completely eliminates the regression with no architectural changes. The fix is a data preparation technique, not a model modification. But the fix was only possible because per-grammar-role decomposition detected the regression and motivated the analysis that identified its cause.

We recommend per-grammar-role decomposition as standard practice for structured output fine-tuning, and presupposition labeling as a principled approach to fields whose values are not always determinable from the input.

\section*{Acknowledgments}

The author used Claude (Anthropic) for code development and manuscript preparation. Experiments were conducted on RunPod GPU instances. The Schema-Guided Dialogue dataset \cite{rastogi2020sgd} is used under CC BY-SA 4.0 license.

\bibliography{references}
\bibliographystyle{plainnat}

\end{document}
